\section{Desenvolvimento do Memora}
\label{sec:desenvolvimento}

Esta seção apresenta os resultados das etapas que antecederam a implementação: a comparação dos simuladores e a escolha da base, o levantamento e a definição dos requisitos, a inspeção de usabilidade e o projeto do protótipo. A exposição relaciona o que foi observado às decisões adotadas para o Memora.

\subsection{Análise comparativa e escolha da ferramenta de base}

A análise considerou ferramentas relacionadas à memória cache que possuíam acesso gratuito ou acadêmico e documentação que permitisse caracterizar seus recursos. Foram incluídos simuladores visuais e ferramentas por linha de comando, para examinar diferentes formas de apresentar estados e resultados. A disponibilidade de código foi um critério para escolher a base a ser modificada, sem excluir da comparação ferramentas gratuitas cujo código público não foi identificado.

